% =============================================================
% Subsection: Model (inside Section 2)
% =============================================================

A representative firm in group $g\in\{\text{Small},\text{Large}\}$ produces output $Y_g$ from capital $K_g$ and effective labor $L_g$ using a Cobb-Douglas technology. Small firms have at most 49 employees, and large firms have 50 or more.
\begin{equation}\label{eq:production}
    Y_g = A\cdot K_g^{\alpha}\cdot L_g^{1-\alpha},
\end{equation}
where $A$ is aggregate productivity and $\alpha$ is the capital share of income. Effective labor $L_g$ combines formal ($F$) and informal ($I$) workers through a CES aggregator,
\begin{equation}\label{eq:ces}
    L_g = \left[\omega L_{F,g}^{\rho} + (1-\omega)L_{I,g}^{\rho}\right]^{1/\rho},\quad \rho=(\sigma_{\text{sub}}-1)/\sigma_{\text{sub}},
\end{equation}
where $\omega$ is the CES weight on formal labor and $\sigma_{\text{sub}}$ is the elasticity of substitution between formal and informal labor. When $\sigma_{\text{sub}}<1$ the two inputs are complements, and raising the cost of formal labor does not trigger a large reallocation toward the informal margin. The two labor inputs are
\[
L_{F,g}=N_{F,g}\sum_b\theta_b h_b e(h_b),\qquad L_{I,g}=\eta_I N_{I,g} h_I e(h_I),
\]
where $N_{F,g}$ and $N_{I,g}$ are the headcounts of formal and informal workers in group $g$ (with $N_{I,g}\equiv N_g-N_{F,g}$ and total group employment $N_g$ fixed), $h_b$ indexes the three contracted-hours bins $b\in\{36,40,44\}$ with shares $\theta_b$ among formal workers, $h_I=44$ is weekly hours for informal workers, and $\eta_I$ is the per-head productivity of an informal worker relative to a formal worker. Per-hour efficiency follows $e(h)=\exp\{-\kappa(h-h^*)^2\}$ above the peak $h^*=40$, where $\kappa$ is the curvature of the long-hours fatigue penalty. Below the peak the paper considers two transparent regimes. The preferred \emph{one-sided fatigue} specification sets $e(h)=1$ for $h\leq h^*$: long workweeks are tiring, but shorter workweeks are not mechanically punished. The conservative \emph{two-sided penalty} specification keeps the same quadratic form below $h^*$: deviations on either side of 40 hours reduce per-hour efficiency. The one-sided specification is informed by \citet{Pencavel2015}, who documents the fatigue leg above the peak but does not characterize efficiency below it. The two-sided specification serves as a conservative upper bound on the required TFP gain.

Each firm chooses the number of formal workers $N_{F,g}$ to maximize profit net of three cost wedges,
\begin{equation}\label{eq:firm}
    Y_g-\tau_g N_{F,g}-\tfrac{1}{2}\pi_m N_{I,g}^2-\tfrac{1}{2}\gamma_{F,g}(N_{F,g}-\bar{N}_{F,g})^2,
\end{equation}
where $\tau_g$ is a linear per-worker cost of formality (CLT/FGTS obligations and size-dependent enforcement), $\pi_m$ is a convex cost of informality (fines, reputational risk, and limited access to formal inputs), and $\gamma_{F,g}$ is a quadratic adjustment cost penalizing deviations from the pre-reform formal headcount $\bar{N}_{F,g}$. The wedge $\tau_g$ is calibrated to match group-specific informality and plays the role of a size-dependent distortion in the spirit of \citet{HsiehKlenow2009} and \citet{RestucciaRogerson2008}, with endogenous-formality microfoundations in \citet{DixCarneiroGoldbergMeghirUlyssea2026}. Cutting the statutory cap from $h_0=44$ hours to $h_1$ reduces effective labor through both the hours bins $h_b$ and the efficiency function $e(h_b)$, and widens the wedge between the marginal product of formal labor and its private cost.

The object of interest is the TFP multiplier that restores aggregate output,
\begin{equation}\label{eq:areq}
    A_{\text{req}}:\quad \sum_g Y_g\!\left(A\,(1+A_{\text{req}}),K_g,h_1\right)=\sum_g Y_g(A,K_g,h_0).
\end{equation}
$A_{\text{req}}$ is the one-off percentage gain in aggregate productivity $A$ that exactly offsets the short-run output loss from moving the cap from $h_0$ to $h_1$. It is an accounting benchmark to compare with Brazil's historical TFP record, not a predicted growth rate.

For welfare, preferences are GHH \citep{GreenwoodHercowitzHuffman1988} with $\sigma_{\text{ghh}}=1$ and Frisch elasticity $1/\nu=0.5$. The compensating variation is
\begin{equation}\label{eq:cv}
    \Delta\text{CV}=\frac{C_1-\psi h_1^{1+\nu}/(1+\nu)}{C_0-\psi h_0^{1+\nu}/(1+\nu)}-1.
\end{equation}
The numerator and denominator are the post- and pre-reform consumption-equivalents of welfare under GHH preferences, computed as consumption net of the monetized labor disutility $\psi h^{1+\nu}/(1+\nu)$. $\Delta$CV is the percentage change in welfare.

The model is static and partial-equilibrium. Capital is predetermined, total employment $N$ is fixed, and general-equilibrium feedbacks (wage-price, monetary, fiscal, cross-sector) are absent. Appendix Table~\ref{tab:omitted_channels} lists the modeled, bracketed, and omitted channels, including compliance, overtime, wage bargaining, labor supply, and sectoral linkages. A textbook user-cost elasticity $\varepsilon_K\in[0.2,0.4]$ \citep{Caballero1999} with capital share $\alpha=0.35$ implies
\begin{equation}\label{eq:pe_bound}
    \Delta A_{\text{req}}^{K}\approx -\alpha\cdot\varepsilon_K\cdot|\Delta L_F/L|\in[-1.0,-1.8]\text{ pp},
\end{equation}
where $\Delta A_{\text{req}}^{K}$ is the downward correction to the short-run $A_{\text{req}}$ once the endogenous response of the capital stock is factored in. The medium-run analog of the headline falls by $1$ to $2$ pp over a $3$ to $5$ year horizon.
